%==============================================================
% LAMPIRAN
%==============================================================

% Halaman standalone judul "LAMPIRAN" di tengah halaman
\clearpage
\vspace*{\fill}
\begin{center}
  {\fontsize{14pt}{17pt}\selectfont\bfseries LAMPIRAN}
\end{center}
\vspace*{\fill}
\addcontentsline{toc}{chapter}{LAMPIRAN}

%--------------------------------------------------------------
% LAMPIRAN 1 (halaman tersendiri)
%--------------------------------------------------------------
\clearpage

\refstepcounter{lampiran}
\addcontentsline{lamp}{lampiran}{\protect\numberline{\thelampiran}Lampiran 1\enspace Rata-rata dan simpangan baku beberapa sifat fisik dan kimia tanah dari 78 contoh tanah di Kebun Percobaan Ciheuleut}
\label{lamp:tanah}

% Caption lampiran: di atas tabel, dengan format "Lampiran N Judul"
\noindent\hangindent=5.6em
\makebox[5.4em][l]{\textbf{Lampiran \thelampiran}}%
Rata-rata dan simpangan baku beberapa sifat fisik dan kimia tanah dari
78 contoh tanah di Kebun Percobaan Ciheuleut

\vspace{4pt}

\centering
\begin{tabular}{lcc}
  \toprule
  Sifat
    & Rata-rata
    & Simpangan baku \\
  \midrule
  Pasir (\%)                        & 47.66 & 23.81 \\
  Lempung (\%)                      & 21.80 & 11.94 \\
  Liat (\%)                         & 30.72 & 18.09 \\
  C-organik (\%)                    & 0.61  & 0.57  \\
  Rapatan isi (mg m$^{-3}$)         & 1.43  & 0.16  \\
  KTK (mek 100 g$^{-1}$ tanah)$^{\mathrm{a}}$
                                    & 18.08 & 17.09 \\
  KAT pada KL (g g$^{-1}$)         & 23.62 & 10.80 \\
  KAT pada TLP (g g$^{-1}$)        & 11.11 & 9.05  \\
  \bottomrule
\end{tabular}

\begin{minipage}{0.85\textwidth}
  \vspace{3pt}
  \footnotesize
  $^{\mathrm{a}}$Banyaknya 70 contoh tanah; KTK: kapasitas tukar kation,
  KAT: kadar air tanah, KL: kapasitas lapang, TLP: titik layu permanen.
\end{minipage}

\clearpage

%--------------------------------------------------------------
% LAMPIRAN 2
%--------------------------------------------------------------

\refstepcounter{lampiran}
\addcontentsline{lamp}{lampiran}{\protect\numberline{\thelampiran}Lampiran 2\enspace Umur, indeks luas daun, dan hasil biji kering jagung yang ditanam pada lima ketinggian tempat}
\label{lamp:jagung}

\noindent\hangindent=5.6em
\makebox[5.4em][l]{\textbf{Lampiran \thelampiran}}%
Umur, indeks luas daun, dan hasil biji kering jagung yang ditanam pada
lima ketinggian tempat

\vspace{4pt}

\noindent
\begin{tabular}{cccc}
  \toprule
  Ketinggian (m dpl) & Umur (hari) & Indeks luas daun & Hasil (ton ha$^{-1}$) \\
  \midrule
  856 & 115 & 3.10 & 5.69 \\
  605 & 106 & 3.09 & 5.43 \\
  400 & 100 & 2.47 & 4.80 \\
  210 &  93 & 2.46 & 4.25 \\
   10 &  88 & 2.12 & 4.03 \\
  \bottomrule
\end{tabular}
